\scanpage{35}
\diagram{\begin{tikzcd}[column sep=small]
& X\arrow[d]\arrow[dl,dashed,"\phi"'] & Y\arrow[dl]\\
S & T\arrow[l,"f"'] & T\times_S T\arrow[l,bend right=14,"\operatorname{pr}_1"']\arrow[l,bend left=14,"\operatorname{pr}_2"]\arrow[ll,bend left=35,"g"']
\end{tikzcd}}
\[
f_*(X/T)=\underline{\operatorname{Hom}}_{T/S}(T/S,X/S)
\]
\[
f_*(X/T)=\underline{\operatorname{Hom}}_S(T,X)\times_{\underline{\operatorname{Hom}}_S(T,T)}S
\]
\[
\underline{\operatorname{Hom}}_{T/S}(X/T,Y/T)=f_*(\underline{\operatorname{Hom}}_T(X,Y))
\]
\[
\underline{\operatorname{Hom}}_S(X,Y)=\phi_*(X\times Y/X)
\]
\[
\underline{\operatorname{Hom}}_S(X,Y)=\underline{\operatorname{Hom}}_{S/S}(X/S,Y/S)
\]

Si $X$ est muni d'une donnée de descente à $S$, et si $f_*(X)$ existe \inserted{[et si on a les conditions d'\unsure{immersité|immersivité}\ldots]}, alors l'objet descendu $X_0$ existe aussi, et on aura
\[
X_0=f_*(X)\times_{g_*(X\times_S T)}f_*(X)
\]

Contre-exemple de Nagata : schéma algébrique simple $X$, \inserted{\unsure{absolument|de dimension 3}}, sur un corps alg. clos $k$, avec groupe $G\simeq\mathbb Z/2$ opérant sans point fixe, tel que $X/G$ n'existe pas. Soit \unsure{$L$} une extension de $k$ sur laquelle $G$ opère \uline{fidèlement}, \deleted{pt fixe} $K=L^G$, faisons opérer $G$ \uline{diagonalement} sur $X\otimes_k L/L$, on a donc sur $X\otimes_k L/L$ une donnée de \uline{descente galoisienne} pour $L/K$. Je dis que \uline{le quotient n'existe pas}, i.e. \uline{la descente n'est pas possible}. En effet, \deleted{si le quotient existait, il existerait aussi par extension des scalaires \inserted{de $K$} à $\unsure{\bar L|L}$,
\[
(X\otimes_k L)\otimes_K\unsure{\bar L|L}=X\otimes_k(L\otimes_K L)=X\otimes_k L\unsure{\longrightarrow|\text{—}}
\]}
% archive foliation: 34
